% BANKS-SOURCE: pdf=111 print=101 kind=section
\subsection{The Higgs model, duality, and the phases of gauge theory}

% BANKS-SOURCE: pdf=111 print=101 kind=prose
Let us recall the Stueckelberg formalism for massive vector fields. We describe
a massive vector $B_\mu$ as a gauge-invariant theory of a Maxwell field
$A_\mu$ coupled to a scalar, with the Lagrangian
% BANKS-SOURCE: pdf=112 print=102 kind=prose
\begin{equation*}
  \mathcal L_{\mathrm{Stueck}}
  =-\frac{1}{4g^2}A_{\mu\nu}^2
  -\frac{\mu^2}{2}(A_\mu-\partial_\mu\theta)^2.
\end{equation*}
Note that, unlike a conventional scalar field, $\theta$ is dimensionless.

The Higgs model is obtained by the simple observation that the Stueckelberg
Lagrangian may be viewed as an approximation to a more conventional Lagrangian
of a charged scalar field coupled to electromagnetism. Simply write
\begin{equation*}
  \mathcal L_{\mathrm{Higgs}}=-\frac{1}{4g^2}A_{\mu\nu}^2
  -\frac{1}{g^2}[|D_\mu\phi|^2+V(\phi^*\phi)].
\end{equation*}
Transforming to radial coordinates for the charged field,
$\phi\equiv\rho e^{\ii\theta}$, we obtain
\begin{equation*}
  \mathcal L_{\mathrm{Higgs}}=-\frac{1}{4g^2}A_{\mu\nu}^2
  -\frac{1}{g^2}[\rho^2(A_\mu-\partial_\mu\theta)^2
  +(\partial_\mu\rho)^2+V(\rho)].
\end{equation*}
If $\rho$ is ``frozen out'' at a value $\rho^2=\phi_0^2=\mu^2/2$, we obtain
the Stueckelberg Lagrangian. In the Higgs model, $\rho$ fluctuates around its
expectation value and quantization of the small fluctuations gives another
particle, called the Higgs boson.

The renormalizable version of this Lagrangian contains two parameters, in
addition to the gauge coupling, $g^2$. The Higgs potential is
\begin{equation*}
  V=\lambda(\phi^*\phi-\phi_0^2)^2.
\end{equation*}
It has a one-parameter set of minima,
\begin{equation*}
  \phi_m=\phi_0e^{\ii\alpha},
\end{equation*}
which are related by gauge transformations (and thus represent a redundant
parametrization of the gauge-invariant minimum $\rho=\phi_0$). Small
fluctuations around a minimum, in the direction of the gauge transformation,
are the degrees of freedom associated with the Stueckelberg field $\theta$ and
are ``eaten'' by the gauge boson $A_\mu$ to form the massive field $B_\mu$.
Fluctuations perpendicular to the trough of minima are Higgs particles. The
gauge-boson mass is
\begin{equation*}
  m_V^2=2\phi_0^2,
\end{equation*}
and the Higgs boson mass
\begin{equation*}
  m_H^2=4\lambda\phi_0^2.
\end{equation*}
The Stueckelberg Lagrangian is a good low-energy approximation if
\begin{equation*}
  \lambda\gg1.
\end{equation*}
The quartic coupling of canonically normalized fields is $g^2\lambda/4$, so
this can be true in a semi-classical regime if $g^2\ll1$.

% BANKS-SOURCE: pdf=113 print=103 kind=prose
In order to quantize the theory, we need to choose a gauge-fixing function $F$
and an associated gauge fermion. The standard covariant choice
$F=\partial_\mu A^\mu$ recommends itself, as does the alternative covariant
choice
\begin{equation*}
  F=(\phi-\phi^*)\phi_0,
\end{equation*}
fixing the gauge freedom in $\phi$ completely.\footnote[7]{Strictly speaking
$\rho(x)=(\phi+\phi^*)(x)$ is positive, but in perturbation theory it is
expanded about the positive value $\phi_0$, and this constraint is fulfilled
automatically.} This, the so-called unitary gauge,\footnote[8]{Also called the
unitarity gauge.} is covariant and contains no unphysical degrees of freedom.
The gauge potential in this gauge is actually equal to the gauge-invariant
field $B_\mu$, because the phase of $\phi$ is frozen. Unfortunately, in
non-abelian theories, Green functions in the unitary gauge do not have a
renormalizable perturbation expansion (although the scattering matrix does --
see below). There is actually a convenient interpolating choice for the gauge
fermion. We take
\begin{equation*}
  \Psi=\frac{1}{g^2}\bar c(\partial_\mu A^\mu
  +\ii\alpha\phi_0(\phi-\phi^*)).
\end{equation*}
This is to be used only with a weight function for the Lagrange multiplier
field
\begin{equation*}
  \delta\mathcal L=\frac{1}{2g^2\alpha}N^2.
\end{equation*}
The propagators in the semi-classical expansion in this gauge contain no mixing
between the scalar and gauge field for any value of $\alpha$.

The Feynman rules for a general, non-abelian, Higgs model in this $R_\alpha$
gauge are given in Appendix D. We will not do any explicit loop computations
in Higgs models in the text, but leave them for the exercises.

\BanksImplicitHook{I-CH08-005}
